% !PRINTED 221
\weberpage{221}

Die Kongruenz (9) besitzt, wenn $k$ die Anzahl der in $n$ aufgehenden, voneinander
verschiedenen Primzahlen ist, $2^k$ inkongruente Lösungen, und dies ist also der Grad
der Gruppe $\mathfrak{B}_0$\footnote{Vgl. \gesperrt{Dirichlet-Dedekind}, Vorlesungen über Zahlentheorie, 4.\ Aufl., § 37.}.

Eine rationale Funktion $\xi$ der Wurzeln \gesperrt{einer Reihe}, etwa der Reihe $R_{1,0}$,
läßt sich nach (1) rational durch \gesperrt{eine dieser} Wurzeln darstellen. Wenn diese Funktion
die Eigenschaft hat, ungeändert zu bleiben, falls diese eine Wurzel durch eine andere
derselben Reihe ersetzt wird, wenn also beispielsweise $\xi$ eine \gesperrt{symmetrische Funktion
der Wurzeln einer Reihe} ist, dann erhält $\xi$ durch Anwendung der Substitutionen von
$\mathfrak{A}$ (oder auch von $\mathfrak{B}$) nur
\weq{10}{\nu = \psi(n)}
verschiedene Werte
\weq{11}{\xi_1, \xi_2, \ldots, \xi_\nu,}
% Source reading \xi_\nu retained: the damaged subscript is fixed by the immediately following count of \nu quantities.
\gesperrt{und wenn diese Werte alle voneinander verschieden sind, so gehört $\xi$ zu der Gruppe
$\mathfrak{A}_{1,0}$ und alle anderen symmetrischen Funktionen der Wurzeln von $R_{1,0}$
sind rational durch $\xi$ darstellbar} (Bd.\ I, § 162).

Die $\nu$ Größen (11) sind die Wurzeln einer irreduzibeln rationalen Gleichung $\nu$ten
Grades, die wir eine \gesperrt{Transformationsgleichung} nennen.

Die Transformationsgleichung hat, wenn in $\xi$ nur rationale Zahlkoeffizienten
vorkommen, selbst rationale Zahlkoeffizienten. \gesperrt{Sie bleibt aber irreduzibel, wenn auch
$n$te Einheitswurzeln oder überhaupt irgend welche Konstanten adjungiert werden}, wie
sich daraus ergibt, daß durch Substitutionen der Gruppe $\mathfrak{B}$ jede Wurzel der
Teilungsgleichung in jede andere, also auch jede Reihe in jede andere Reihe übergeführt
werden kann.

Durch die Adjunktion einer Wurzel $\xi$ der Transformationsgleichung, etwa der zur
Gruppe $\mathfrak{A}_{1,0}$ gehörigen, reduziert sich die Gruppe der Teilungsgleichung
auf $\mathfrak{A}_{1,0}$, die letztere Gruppe ist aber nicht mehr transitiv, sondern
vertauscht die $\varphi(n)$ Wurzeln $x_{h,0}$, worin $h$ ein vollständiges System
inkongruenter, zu $n$ teilerfremder Zahlen durchläuft, untereinander. Die
Teilungsgleichung wird also reduzibel und hat einen Faktor $\varphi(n)$ten Grades,
%CONT>>
